\subsection{Aufgabe 1 - Richtmomentbestimmung}
Durch das Auftragen von \(\varphi(M)\) die Steigung \(a\) abgelesen werden, daraus ergibt sich \(D\) nach \cref{M_linear} durch:
\begin{equation}
    D=\frac{1}{a}=\frac{M}{\varphi}
    \label{D_a}
\end{equation}
Das Auftragen erfordert die Berechnung von \(M\). Hierfür gilt nach \cref{M}:
\begin{align}
  M=m\,g\,r&&  \Delta M=\sqrt{\left(g r \Delta m\right)^2+\left(g m \Delta r\right)^2}
\end{align}
wobei \(m\) die fehlerbehaftete Masse der Gewichtstellermassen ist und \(r=\tfrac{1}{2}d\) für \(d\) als gemessenen Durchmesser der Scheibe.
Mit der Fehlerformel, sowie den Fehlerwerten von \(\Delta m=\qty{e-4}{\kilogram}\), \(\Delta r=\qty{0.0025}{\cm}\) (wurde aus \(\Delta d=\qty{0.005}{\cm}\) fortgepflanzt), ergibt sich die Tabelle unter Nutzen von \(g=\qty{9.81}{\meter\per\second\squared}\):
\begin{table}[htbp]
  \centering
  \caption{berechnete Drehmomente für \cref{fig:Aufgabe1_4.pdf}}
  \begin{tabularx}{0.8\textwidth}{ZZZZ}
    \toprule
    \text{Masse }m [\unit{\kg}] & \varphi [\unit{\degree}] & M [\qty{e-3}{\Nm}]& \Delta M [\qty{e-3}{\Nm}]\\ 
    \midrule            
          0.05            & 58\pm2  & 23.30 & 0.14                    \\
          0.1             & 120\pm2 & 47.60 & 0.25                    \\
          0.15            & 183\pm2 & 69.9  & 0.4                     \\
          0.2             & 242\pm2 & 93.2  & 0.5                     \\
          0.25            & 306\pm2 & 116.5 & 0.7                     \\
          0.3             & 368\pm2 & 139.8 & 0.8                     \\
    \bottomrule
  \end{tabularx}
\label{table:Drehmomente}  
\end{table}

Als Ergebnis für \(a\) wurde grafisch in \cref{fig:Aufgabe1_4.pdf} ermittelt: \(\tcbhighmath[colback=white]{a=(2.62\pm0.24)\qty{e3}{\degree\per\Nm}}\)\\
Damit berechnet sich durch \cref{D_a} und \(\Delta D=\tfrac{\Delta a}{a^2}\):
\begin{rbl}
D_1=(0.38\pm0.04)\qty{0.38(4)e-3}{\Nm\per\degree}=\qty{0.0218(23)}{\Nm\per\radian}
\end{rbl}

\xfigpdf{htbp}{page=1, width=0.9\textwidth}{Aufgabe1_4.pdf}{Einzeichnen von \(\varphi(M)\) zur grafischen Bestimmung von \(a\)}{Einzeichnen von \(\varphi(M)\) zur grafischen Bestimmung von \(a\)}
\clearpage

\subsection[\texorpdfstring{Aufgabe 2 - \(D\) durch Nutzen der runden Messingscheibe}{Aufgabe 2 - D durch Nutzen der runden Messingscheibe}]{Aufgabe 2 - \(\boldsymbol{D}\) durch Nutzen der runden Messingscheibe}
Für einen Zylinder wie die Scheibe mit Radius \(r\) und Masse \(m\) gilt:
\begin{equation}
  J_S=\frac{1}{2}m\,r^2
\end{equation}
Weiß man nun, dass ein Körper, der aus zwei Subkörpern besteht, deren Schwerpunkt übereinander liegt (Drehtisch mit \(J_T\) und runde Messingscheibe mit \(J_S\)), ein Trägheitsmoment von \(J=J_T+J_S\) besitzt, und die Periodendauern \(T_1\) (nur der Drehtisch montiert) und \(T_2\) (Drehtisch+Scheibe) bekannt sind, dann können die Formeln:
\begin{align}
  T_1=2\pi\sqrt{\frac{J_T}{D}}&&
  T_2=2\pi\sqrt{\frac{J_T+J_S}{D}}
\end{align}
durch einfaches Subtrahieren miteinander verknüpft werden, zu:
\begin{align}
  D=&\frac{4 \pi^{2} J_{s}}{T_{2}^{2} - T_{1}^{2}}
  \label{J_Platte}\\
  =& \frac{2 \pi^{2} m r^{2}}{T_{2}^{2} - T_{1}^{2}}
  \label{D_T}
\end{align} 
Man beachte, dass das Trägheitsmoment \(J_T\) des Tisches sich heraussubtrahiert hat. Der Fehler berechnet sich mit:
\begin{equation}
  \begin{aligned}
 \Delta D^2=&\left(\frac{2{\pi}^2 r^2}{-{T_1}^2+{T_2}^2} \Delta m\right)^2+\left(\frac{4{\pi}^2 m r}{-{T_1}^2+{T_2}^2} \Delta r\right)^2\\
 +&\left(\frac{4{\pi}^2 T_1 m r^2}{\left(-{T_1}^2+{T_2}^2\right)^2} \Delta T_1\right)^2+\left(\frac{-4{\pi}^2 T_2 m r^2}{\left(-{T_1}^2+{T_2}^2\right)^2} \Delta T_2\right)^2
  \end{aligned}
\end{equation}
Es berechnet sich:
\begin{rb}
D_2=\qty{0.0225(4)}{\Nm\per\radian}
\end{rb}
Dies stimmt innerhalb der Fehlergrenzen mit \(D_1\) überein.
\newpage

\subsection[\texorpdfstring{\(J\) der runden Messingplatte}{J der runden Messingplatte}]{\(\boldsymbol{J}\) der runden Messinglatte}
Weil es mich interessiert hat, habe ich einmal mit:
\begin{equation}
  J=\frac{1}{2}m\,r_\perp^2
\end{equation}
\begin{equation}
  \Delta J^2=\left(\frac{r^2}{2} \Delta m\right)^2+\left(m r \Delta r\right)^2
\end{equation}
das Trägheitsmoment der runden Scheibe ausgerechnet. Es ergibt sich:
\[J_{\text{allgemein}}=\qty{0.9860(18)e-3}{\kg\meter\squared}\]
Mit der Bestimmung über den Weg der folgenden Aufgabe ergibt sich:
\[J_{\text{Umweg}}=\qty{0.981(23)e-3}{\kg\meter\squared}\]
Das ist eine Abweichung von \(0.22\sigma\). 

\subsection[\texorpdfstring{Aufgabe 3 - \(J_P\) der umförmigen Platte}{Aufgabe 3 - J P der unförmigen Platte}]{Aufgabe 3 - \(\boldsymbol{J_P}\) der unförmigen Platte}
Mithilfe der \cref{J_Platte} und der Periodendauer von \(T_3\) der unförmigen Scheibe lässt sich bei bekanntem \(D\) berechnen: 
\begin{equation}
  J_P=\frac{D(T_3^2-T_1^2)}{4\pi^2}
  \label{J_P}
\end{equation}
\(D\) und die Periodendauern sind fehlerbehaftet: 
\begin{equation}
  \Delta J_P^2=\left(\frac{-{T_1}^2+{T_3}^2}{4{\pi}^2} \Delta D\right)^2+\left(\frac{-D T_1}{2{\pi}^2} \Delta T_1\right)^2+\left(\frac{D T_3}{2{\pi}^2} \Delta T_3\right)^2
\end{equation}
Da wir zwei Ergebnisse für \(D\) haben berechnet man den gewichteten Mittelwert \autocite{Mittelwert}. Dafür werden \(D_1\),\(_2\) mit Gewichten \(g_m\) (\(m=(1,2))\) versehen, die sich wie folgt berechnen:
\begin{equation}
g_m=\frac{1}{\Delta D_m^2}
\end{equation}
das gewichtete Mittel berechnet sich mit:
\begin{align}
\bar{D}=\frac{\sum_{m=1}^{2} g_mD_ m}{\sum_{m=1}^2 g_m}&&\Delta \bar{D}=\sqrt{\frac{1}{\sum_{m=1}^{2} g_m}}
\end{align}
Dann ergibt sich:
\(\tcbhighmath[colback=white]{\bar{D}=\qty{0.0224(4)}{\Nm\per\radian}}\).\\
Nun nutzen wir noch unser Ergebnis für \(T_3=\qty{2.33}{\s}\) mit \(\Delta T_3=\frac{\qty{0.2}{\s}}{20}=\qty{0.01}{\s}\), dies ist der Reaktionszeitfehler über 20 Schwingungen, und das vorhin schon benutzte \(T_1=\qty{1.200(3)}{\s}\). Es ergibt sich:
\begin{rb}
J_{P_S}=\qty{2.26(5)e-3}{\kg\meter\squared}
\end{rb}

\subsection{Steinerscher Satz}
Wir berechnen für unterschiedliche Abstände \(a_i\) zur Schwerpunktachse mithilfe der \cref{J_P} und den gemessenen Periodendauern die Trägheitsmomente \(J_{P_i}\) und einmal mithilfe von Steiner \cref{Steiner}, hierfür wird das eben berechnete \(J_{P_S}\) genutzt sowie die gewogene Masse \(m=\qty{0.695(1)}{\kg}\) und der Abstandsfehler \(\Delta a=\qty{0.5}{\mm}\) (wurde nur mit einem Geodreieck abgemessen).
Der Fehler des Steinerschen Moments ergibt sich durch:
\begin{equation}
  \Delta J_{i,\,\text{Steiner}}^2=\left(\Delta J_{P_S}\right)^2+\left(a_i^2 \Delta m\right)^2+\left(2a_i\, m \Delta a_i\right)^2
\end{equation}
\(J_{\text{gemessen}}\) wurde wie oben analog mit \cref{J_P} und der zugehörigen Fehlergleichung berechnet. In dieser Tabelle sind die berechneten Werte zu finden. 

\begin{table}[htbp]
  \centering
    \caption{Trägheitsmomente der unförmigen Scheibe auf unterschiedliche Weise berechnet}
  \begin{tabularx}{\textwidth}{xZZZS[table-format=1.2]}
    \toprule
      a [\unit{\cm}]  & a^2 [\unit{\cm\squared}] & J_{\text{gemessen}}[\qty{e-3}{\kg\meter\squared}]  & J_{\text{Steiner}} [\qty{e-3}{\kg\meter\squared}]  &  {\(S[\sigma]\)} \\
      \midrule
      0,5 & 0.25    & (2.32\pm0.05) & (2.28\pm0.05) & 0.6\\
      1   & 1       & (2.37\pm0.05) & (2.33\pm0.05) & 0.6 \\
      1,5 & 2.25    & (2.51\pm0.05) & (2.42\pm0.05) & 1.3\\
      2   & 4       & (2.59\pm0.06) & (2.54\pm0.05) & 0.7\\
      2,5 & 6.25    & (2.81\pm0.06) & (2.69\pm0.05) & 1.5\\
    \bottomrule
  \end{tabularx}
\label{table:J_P_I}
\end{table}

Die Werte \(J(a^2)\) mit entsprechenden Fehlerbalken (auch für \(a\)) finden sich in \cref{fig:Aufgabe4_1.pdf}.

\xfigpdf{htbp}{page=1, width=0.9\textwidth}{Aufgabe4_1.pdf}{}{Eintragen der richtig berechneten Werte}

Hier sind die alten, falschen Werte, auf denen \cref{fig:Aufgabe4.pdf} basiert:

\begin{table}[htbp]
  \centering
  \caption{Trägheitsmomente der unförmigen Scheibe falsch berechnet}
  \begin{tabularx}{\textwidth}{xZZZS[table-format=1.2]}
    \toprule
      a [\unit{\cm}]  & a^2 [\unit{\cm\squared}] & J_{\text{gemessen}}[\qty{e-3}{\kg\meter\squared}]  & J_{\text{Steiner}} [\qty{e-3}{\kg\meter\squared}]  & {\(S[\sigma]\)}  \\
      \midrule
      0.5 & 0.25    & (2.11\pm0.05) & (2.08\pm0.05) & 0.5\\
      1   & 1       & (2.16\pm0.05) & (2.13\pm0.05) & 0.5 \\
      1.5 & 2.25    & (2.28\pm0.06) & (2.22\pm0.05) & 0.8\\
      2   & 4       & (2.36\pm0.06) & (2.34\pm0.05) & 0.25\\
      2.5 & 6.25    & (2.56\pm0.06) & (2.49\pm0.05) & 0.9\\
    \bottomrule
  \end{tabularx}
\label{table:J_P_I_falsch}
\end{table}

\xfigpdf{htbp}{page=1, width=0.9\textwidth}{Aufgabe4.pdf}{}{Eintragen der falsch in Tabelle \cref{table:J_P_I_falsch} berechneten Werte}

% \begin{table}[htbp]
%   \centering
%   \caption{}
%   \begin{tabularx}{\textwidth}{ZZZZ}
%     \toprule

%       \midrule

%     \bottomrule
%   \end{tabularx}
% \label{table:}  
% \end{table}
